\documentclass[11pt,a4paper]{article}
\usepackage[T1]{fontenc}
\usepackage[utf8]{inputenc}
\usepackage{amsmath,amssymb,amsthm}
\usepackage[margin=1in]{geometry}
\usepackage[colorlinks=true,linkcolor=blue,urlcolor=blue]{hyperref}
\theoremstyle{plain}
\newtheorem{theorem}{Theorem}
\newtheorem{lemma}[theorem]{Lemma}
\newtheorem{proposition}[theorem]{Proposition}
\newtheorem{corollary}[theorem]{Corollary}
\theoremstyle{definition}
\newtheorem{remark}[theorem]{Remark}

\title{The empirical recurrence for OEIS A210272, proved by transfer matrix}
\author{Adrian Perez Fontelles\\ \small Independent researcher}
\date{31 August 2026}
\begin{document}
\maketitle

\begin{abstract}
OEIS A210272 counts $(n+1)\times5$ arrays over $\{0,\dots,2\}$ subject to a condition imposed on
every $2\times2$ subblock, and carries an empirical recurrence of order $69$ contributed
by R.~H.~Hardin. It is true, and it is not an empirical matter at all. The condition is local
to each subblock, so the rows of an admissible array are the vertices of a finite digraph
on $S=243$ vertices whose edges are the admissible consecutive pairs, and the count is a
number of walks: $a(n)=\tfrac1{2}\mathbf{1}^{\!\top}M^{n}\mathbf{1}$. Writing $q$
for the characteristic polynomial of the conjectured recurrence, the recurrence holds for all
$n$ exactly when $\mathbf{1}^{\!\top}M^{j}q(M)\mathbf{1}=0$ for every $j\ge0$, and the
Cayley--Hamilton theorem bounds that to $j<S$. It is therefore a finite computation in exact
integer arithmetic, and it comes out zero.
\end{abstract}

\noindent\small 2020 Mathematics Subject Classification. 05A15, 05B45, 15B36.\normalsize

\section{The sequence and the conjecture}

OEIS A210272 is ``Half the number of (n+1)X5 0..2 arrays with every 2X2 subblock having two or three distinct clockwise edge differences.''. It has offset $1$ and begins
\[
16620, 2307956, 321050788, 44678808183, 6218088462292, 865403655856588, 120443009719061218, 16762728653468925176,\ \dots
\]
The entry carries the following comment:
\begin{quote}\small
Empirical: a(n) = 142*a(n-1) +1390*a(n-2) -255743*a(n-3) -110496*a(n-4) +169565270*a(n-5) -665183207*a(n-6) -52878902879*a(n-7) +395903765675*a(n-8) +8133200677492*a(n-9) -92576230139190*a(n-10) -542663156454254*a(n-11) +10497292789475974*a(n-12) +2631885281757262*a(n-13) -646106145808214387*a(n-14) +1826441384938919701*a(n-15) +22362684125873346283*a(n-16) -123476318341203413496*a(n-17) -395337011643732047682*a(n-18) +4142092565872768583256*a(n-19) +605782653333712666179*a(n-20) -83518787705085031195627*a(n-21) +136634888376530178952716*a(n-22) +1040134849078032981074015*a(n-23) -3402074914964904960822357*a(n-24) -7239465059314247652144549*a(n-25) +44377737415311976751838594*a(n-26) +9840741933860801981979443*a(n-27) -361590798269938925507784670*a(n-28) +321818684009046992762532501*a(n-29) +1883576717146322164819296451*a(n-30) -3479890138790632205298850558*a(n-31) -5878393363654009470130124838*a(n-32) +19078667220689449803911834552*a(n-33) +7392854341921657854446418742*a(n-34) -65742784478773744652167403362*a(n-35) +19627223671361316822942496555*a(n-36) +149839903328239191224944472173*a(n-37) -117630629870435594695245179193*a(n-38) -226767728498105405443631444233*a(n-39) +284875215327416958030542435899*a(n-40) +219032259251618705695478104114*a(n-41) -424657211702672231000865863606*a(n-42) -114298568153964486240525208337*a(n-43) +426463674134772011008406591220*a(n-44) -1465111296208263054958667449*a(n-45) -298361702108034534055790446739*a(n-46) +50042360589416264853656878477*a(n-47) +147499557244192463582697992372*a(n-48) -39645514435604456524417564586*a(n-49) -51787606560690998053903082159*a(n-50) +16872384336675821704842993456*a(n-51) +12907214080006987000746559571*a(n-52) -4471894138938904280092501035*a(n-53) -2270649243797947159981070488*a(n-54) +757433256288002048852712547*a(n-55) +277898940359626427914691228*a(n-56) -80741549607760164976554158*a(n-57) -22888243320669979845951670*a(n-58) +5181436657383846894790730*a(n-59) +1185029940722673827637793*a(n-60) -185251468573962367963557*a(n-61) -33904216065585670936623*a(n-62) +3275032115968437164460*a(n-63) +413089312994017424066*a(n-64) -25348341919215966472*a(n-65) -1274153174398735370*a(n-66) +12493913728073620*a(n-67) +443471547890804*a(n-68) -2001148574736*a(n-69)
\end{quote}
As of the ``Last modified'' line on the live entry (Oct 03 2025 15:51:33, revision 9) the statement is
still recorded as empirical, and nothing on the entry records it as settled either way.

Written out, the claim is
\begin{equation}\label{eq:conj}
a(n)\;=\;142\,a(n-1) + 1390\,a(n-2) - 255743\,a(n-3) - 110496\,a(n-4) + 169565270\,a(n-5) - 665183207\,a(n-6) - 52878902879\,a(n-7) + 395903765675\,a(n-8) + 8133200677492\,a(n-9) - 92576230139190\,a(n-10) - 542663156454254\,a(n-11) + 10497292789475974\,a(n-12) + 2631885281757262\,a(n-13) - 646106145808214387\,a(n-14) + 1826441384938919701\,a(n-15) + 22362684125873346283\,a(n-16) - 123476318341203413496\,a(n-17) - 395337011643732047682\,a(n-18) + 4142092565872768583256\,a(n-19) + 605782653333712666179\,a(n-20) - 83518787705085031195627\,a(n-21) + 136634888376530178952716\,a(n-22) + 1040134849078032981074015\,a(n-23) - 3402074914964904960822357\,a(n-24) - 7239465059314247652144549\,a(n-25) + 44377737415311976751838594\,a(n-26) + 9840741933860801981979443\,a(n-27) - 361590798269938925507784670\,a(n-28) + 321818684009046992762532501\,a(n-29) + 1883576717146322164819296451\,a(n-30) - 3479890138790632205298850558\,a(n-31) - 5878393363654009470130124838\,a(n-32) + 19078667220689449803911834552\,a(n-33) + 7392854341921657854446418742\,a(n-34) - 65742784478773744652167403362\,a(n-35) + 19627223671361316822942496555\,a(n-36) + 149839903328239191224944472173\,a(n-37) - 117630629870435594695245179193\,a(n-38) - 226767728498105405443631444233\,a(n-39) + 284875215327416958030542435899\,a(n-40) + 219032259251618705695478104114\,a(n-41) - 424657211702672231000865863606\,a(n-42) - 114298568153964486240525208337\,a(n-43) + 426463674134772011008406591220\,a(n-44) - 1465111296208263054958667449\,a(n-45) - 298361702108034534055790446739\,a(n-46) + 50042360589416264853656878477\,a(n-47) + 147499557244192463582697992372\,a(n-48) - 39645514435604456524417564586\,a(n-49) - 51787606560690998053903082159\,a(n-50) + 16872384336675821704842993456\,a(n-51) + 12907214080006987000746559571\,a(n-52) - 4471894138938904280092501035\,a(n-53) - 2270649243797947159981070488\,a(n-54) + 757433256288002048852712547\,a(n-55) + 277898940359626427914691228\,a(n-56) - 80741549607760164976554158\,a(n-57) - 22888243320669979845951670\,a(n-58) + 5181436657383846894790730\,a(n-59) + 1185029940722673827637793\,a(n-60) - 185251468573962367963557\,a(n-61) - 33904216065585670936623\,a(n-62) + 3275032115968437164460\,a(n-63) + 413089312994017424066\,a(n-64) - 25348341919215966472\,a(n-65) - 1274153174398735370\,a(n-66) + 12493913728073620\,a(n-67) + 443471547890804\,a(n-68) - 2001148574736\,a(n-69)
\end{equation}
for all $n>68$.

\section{The condition is local, so the rows form a digraph}

Write an array of the shape in question as its list of rows
$r^{(1)},r^{(2)},\dots$, each an element of
$V=\{0,1,\dots,2\}^{5}$, so $|V|=S=243$. A $2\times2$ subblock of the array
occupies two consecutive rows and two consecutive columns; if the two rows are $r$
and $s$ (in that order) and the two columns are numbered $j,j+1$,
the subblock is
\[
\begin{pmatrix}\alpha&\beta\\\gamma&\delta\end{pmatrix}=\begin{pmatrix}r_j&r_{j+1}\\s_j&s_{j+1}\end{pmatrix}.
\]
The entry's requirement on that subblock is
\begin{equation}\label{eq:loc}
\#\{\beta-\alpha,\ \delta-\beta,\ \gamma-\delta,\ \alpha-\gamma\}\in\{2,3\},
\end{equation}
a condition on the four entries alone.

For $r,s\in V$ declare $r\to s$ an edge of a digraph $G$ when, for every
$1\le j\le 5-1$,
\begin{itemize}
\item the four entries above satisfy \eqref{eq:loc};

\end{itemize}
Let $M\in\{0,1\}^{S\times S}$ be the adjacency matrix of $G$ and $\mathbf{1}$ the
all-ones vector.

\begin{lemma}\label{lem:walk}
The arrays counted by the entry with $L$ rows are exactly the walks
$r^{(1)}\to r^{(2)}\to\cdots\to r^{(L)}$ of length $L-1$ in $G$. Consequently
\[
a(n)\;=\;\tfrac1{2}\mathbf{1}^{\!\top}M^{n}\mathbf{1} .
\]
\end{lemma}

\begin{proof}
An array with $L$ rows and $5$ columns is precisely a sequence
$r^{(1)},\dots,r^{(L)}$ of elements of $V$; conversely any such sequence is an array.
Every $2\times2$ subblock lies in two consecutive rows $r^{(t)},r^{(t+1)}$ and two
consecutive columns, so the array satisfies the entry's condition on all of its subblocks
if and only if each consecutive pair $\bigl(r^{(t)},r^{(t+1)}\bigr)$ satisfies it for
every column pair --- which is exactly the edge relation of $G$. The number of walks of
length $L-1$, summed over all starting and ending vertices, is
$\mathbf{1}^{\!\top}M^{L-1}\mathbf{1}$, since $(M^{k})_{r,s}$ counts the walks of
length $k$ from $r$ to $s$. The entry's index $n$ corresponds to $L=n+1$ rows. The entry counts $\tfrac1{2}$ of those arrays, a constant factor that passes through every step below unchanged.
\end{proof}

\section{The criterion}

Let $q(t)=t^{69}-\sum_i c_i\,t^{69-i}$ be the characteristic polynomial of
\eqref{eq:conj}, with the $c_i$ read off that equation.

\begin{lemma}\label{lem:crit}
For every $n$ with $n+1-1\ge69$,
\[
a(n)-\sum_i c_i\,a(n-i)\;=\;\tfrac1{2}\mathbf{1}^{\!\top}M^{\,n+1-1-69}\,q(M)\,\mathbf{1} .
\]
Hence \eqref{eq:conj} holds from that point on if and only if
$\mathbf{1}^{\!\top}M^{j}q(M)\mathbf{1}=0$ for every $j\ge0$; and it is enough to check
$j=0,1,\dots,S-1$.
\end{lemma}

\begin{proof}
By Lemma~\ref{lem:walk}, $a(n-i)=\tfrac1{2}\mathbf{1}^{\!\top}M^{\,n+1-1-i}\mathbf{1}$, so
\[
a(n)-\sum_i c_i a(n-i)
=\tfrac1{2}\mathbf{1}^{\!\top}\Bigl(M^{m}-\sum_i c_i M^{m-i}\Bigr)\mathbf{1}
=\tfrac1{2}\mathbf{1}^{\!\top}M^{\,m-69}\,q(M)\,\mathbf{1}
\]
with $m=n+1-1$, which is the stated identity; the scalar $\tfrac1{2}$ never vanishes, so
it does not affect whether the left side is zero. The equivalence follows by letting
$m-69=j$ run over $\ge0$. For the last clause put $w=q(M)\mathbf{1}$. By the
Cayley--Hamilton theorem $M^{S}$ is an integer combination of $I,M,\dots,M^{S-1}$, so
$\mathbf{1}^{\!\top}M^{j}w$ for $j\ge S$ is the corresponding combination of the earlier
values; if those all vanish, so do all the rest.
\end{proof}

\section{The computation}

\begin{theorem}
The recurrence \eqref{eq:conj} holds for every $n>68$, which is the statement of
Section~1.
\end{theorem}

\begin{proof}
The digraph was constructed from \eqref{eq:loc} by a depth-first scan across the $5$
columns: the edge condition is a conjunction of one constraint per consecutive column
pair, so the successors of a row $r$ are enumerated column by column without testing all
$S^2$ pairs. The vector $w=q(M)\mathbf{1}\in\mathbb{Z}^{S}$ was then formed by $69$
matrix--vector products in exact integer arithmetic, and $\mathbf{1}^{\!\top}M^{j}w$ was
evaluated for $j=0,1,\dots$, again exactly. Every one of those integers vanishes from the
index corresponding to $n=68+1$ onwards, and the run continued until $w$ itself was the
zero vector, which settles every larger $j$ at once. By Lemma~\ref{lem:crit} the recurrence
holds for every $n>68$.
\end{proof}

\begin{remark}
The argument gives more than the recurrence: it shows the sequence is $C$-finite with
characteristic polynomial dividing that of $M$, so it has a rational generating function
whose denominator divides $\det(I-xM)$. The conjectured recurrence is one consequence; the
minimal one is a divisor of $q$.
\end{remark}

\section{Verification}

Three checks, each independent of the algebra above.

First, the digraph was built directly from the entry's stated condition and the walk counts
$\tfrac1{2}\mathbf{1}^{\!\top}M^{L-1}\mathbf{1}$ compared against the entry's DATA at
every one of the $11$ published indices; the values agree exactly. That is what ties
the digraph to the sequence: a model that reproduced only some of the published terms would
be the wrong model, and would have been rejected. In particular it is what rules out a
misreading of the entry's wording, since a different reading of the condition gives a
different digraph and different counts.

Second, the recurrence \eqref{eq:conj} was evaluated on those published terms in exact
integer arithmetic, with no matrices involved, and vanishes wherever it applies.

Third, the annihilation test of Lemma~\ref{lem:crit} was run in exact integer arithmetic
until the working vector was identically zero, not to a sampled prefix, and returned zero
throughout.

\begin{thebibliography}{9}
\bibitem{oeis} The OEIS Foundation, \emph{The On-Line Encyclopedia of Integer
Sequences}, \url{https://oeis.org}, 2026. Sequence A210272.
\bibitem{stanley} R.~P.~Stanley, \emph{Enumerative Combinatorics}, Volume 1, 2nd ed.,
Cambridge University Press, 2011. (Transfer-matrix method, Section 4.7.)
\bibitem{fs} P.~Flajolet and R.~Sedgewick, \emph{Analytic Combinatorics}, Cambridge
University Press, 2009.
\bibitem{horn} R.~A.~Horn and C.~R.~Johnson, \emph{Matrix Analysis}, 2nd ed., Cambridge
University Press, 2013.
\end{thebibliography}

\end{document}
